\chapter{Asking Millionaires Who They're Voting For President}

This chapter follows a School of Hard Knocks interview, curated by LazyingArt LLC, that asks two questions at once and keeps switching the light between them. One question is practical: how was wealth built, preserved, lost, or rebuilt? The other is political: how do wealthy operators say they will vote in the 2024 presidential election? The lecture does not begin with setup. It begins with consequences: a lawyer who used to be a millionaire, a scatter of money figures in the millions, and the election question already in the air. Only after that splice does the host reset us in Austin, Texas. We will keep close to that rhythm, because the chapter works best when the business doctrine is allowed to emerge in the same order as the conversation that produced it.

\section{Cold Open, Austin, and the Two Questions}

The opening montage spends numbers before it spends explanations. We hear rough figures such as \(\$6\text{--}7\) million, \(\$8\) million, \(\$17.5\) million, and \(\$50\) million before we are told clearly whose figures these are or what sort of money they represent. That choice matters. The lecture wants the listener to feel consequence first and reconstruction second.

Only then does the host formalize the scene. We are in Austin, Texas, introduced not merely as a place but as a concentration of wealth:
\begin{equation}
N_{\text{millionaires}} \approx 40{,}000,\qquad
N_{\text{billionaires}} > 10.
\end{equation}
The point of this count is not demography for its own sake. It gives the host permission to ask the same pair of questions repeatedly and expect that the answers may amount to more than casual opinion. From here forward, each interview will be asked to carry a double load:
\begin{itemize}
\item first, reveal a mechanism of wealth, loss, negotiation, or recovery;
\item second, reveal a voting preference or at least a political inclination.
\end{itemize}

We should therefore keep our attention on the lecture's deeper spine. It is not electoral ideology. It is leverage, failure, realization, timing, market pricing, and the conditions under which skill survives after capital has vanished.

\section{Buy Dirt, Fail First, and Negotiate Without Desperation}

The first full interview begins with a wealthy New York builder-developer. The host starts with a broke-versus-wealthy prompt, then line of business, then reputation, and only afterward lets the operating doctrine appear. The speaker identifies himself as a builder-developer, says he worked in both commercial and residential real estate, and compresses a generational lesson into two words: buy dirt.

\paragraph{Anecdote.}
When asked for the best transaction he ever did, he resists the usual trophy-deal answer. Instead he points to streets named after his children. The move is important. In this lecture, wealth is not described only by yearly scale. It is also described by what stays behind in the world.

\paragraph{Claim.}
Asked for the secret to creating wealth in real estate, the answer is failure. One must fail before becoming hungry enough to make. This is not motivational padding. The speaker ties it to an historical contrast: older real-estate conditions offered forms of lift that younger operators no longer inherit so easily.

\paragraph{Mechanism.}
The historical mechanism is described by the phrase ``buy stuff subject to,'' followed by the claim that the property could gain value even before closing. We should keep that cautiously schematic:
\begin{equation}
\text{buy subject to} \;\Rightarrow\; \text{asset appreciates before closing}.
\end{equation}
The interview does not define the legal structure in detail, and we should not pretend that it does. What matters is the speaker's economic memory: there was once a way of controlling an asset early enough to benefit from appreciation before final completion, and he believes that window is narrower now.

The host then narrows from market history to table behavior. What should one do when sitting across from a heavy hitter? The answer is remarkably spare. Do not impulse the deal. Let the other side feel that you do not need it. Above all, do not be desperate. The host's own recap, which comes immediately after the interview, is the correct compression:
\begin{equation}
\text{desperation} \Rightarrow \text{loss of leverage},
\qquad
\text{willingness to walk away} \Rightarrow \text{bargaining power}.
\end{equation}
The speaker makes the same point numerically through a bank-rate anecdote:
\begin{equation}
r_{\text{ask}} = 6\%,\qquad r_{\text{counter}} = 4\%.
\end{equation}
He insists that one should not simply accept the quoted rate, but should counter and then begin walking away if necessary. The underlying model is not a theorem of credit markets. It is a theorem of bargaining stance.

The host also asks for scale and receives a rough, slightly blurred one-year figure:
\begin{equation}
Y_{\text{RE,rough}} \approx \$6\text{--}7\,\text{million/year}.
\end{equation}
The transcript does not cleanly distinguish personal income from business revenue here, and the speaker further muddies the category by saying the business is himself and that he still owns many companies. We should preserve the roughness rather than false precision.

\subsection{Question \& Answer}

Why does willingness to walk away create leverage?

Because the walk-away option changes which side must absorb the pressure of delay. The speaker's financing anecdote is useful precisely because it stays small enough to inspect. If the bank begins at \(6\%\) and the counter is \(4\%\), then the movement is
\begin{align}
\Delta r &= r_{\text{ask}} - r_{\text{counter}} = 2\%,\\
\frac{\Delta r}{r_{\text{ask}}} &= \frac{2}{6} \approx 0.33.
\end{align}
So the anecdotal concession is about one third of the initial ask. We should not mistake this for a universal percentage. The real point is structural. If one side must close now, its feasible set is narrow. If one side can leave, the other side must ask whether losing the deal is worse than improving the terms. This is why the host's recap is stronger than it first sounds: the most desperate person has the least leverage because desperation collapses the set of acceptable next moves.

The political answer arrives only after the business lesson is extracted. The speaker says he will vote for Trump and frames that choice in businessman-versus-businessman terms. In the architecture of the lecture, the vote is the second question, not the first.

\section{Marriage, Paper Wealth, and the Arithmetic of Not Selling}

The next interview changes the register of loss. We move from negotiation at the table to household structure, divorce economics, and then to the harder arithmetic of unrealized wealth. The host again begins with the millionaire question, and the speaker replies by denying current millionaire status while immediately attributing the difference to an ex-wife who has most of the money.

\paragraph{Anecdote.}
The speaker is a lawyer. He says several million dollars were lost in divorce, that the former spouse retained the house, and that in a later marriage he preferred separation to divorce because he did not want retirement wealth further divided away. The host then summarizes the segment with a line that would be crude if it were not so clearly the lecture's intended lesson: marriage is a business.

\paragraph{Claim.}
We should sharpen that claim without overstating it. The interview is not a treatise on family law. It is an argument that household structure enters capital preservation. In the speaker's own language, if one is building wealth, then a prenuptial agreement belongs to the architecture of protection.

\paragraph{Mechanism.}
The deeper mechanism appears when the host asks for the most money made in a single year. The speaker corrects the category immediately: it was paper money, and he lost it because he did not convert it to real money soon enough. That is the mathematical center of the segment:
\begin{equation}
V_{\text{paper}} \neq V_{\text{realized}}
\qquad
\text{when sale constraints or lockups bind.}
\end{equation}
The interview then supplies the values and the timing:
\begin{align}
V_{\text{peak}} &\approx \$17.5\,\text{million},&
\Delta t_{\text{lockup}} &= 6\ \text{months},\\
V_{\text{post-crash}} &\approx \$750{,}000,&
V_{\text{sale}} &\approx \$350{,}000.
\end{align}
In one line, the story is
\begin{equation}
\$17.5\text{M} \longrightarrow \$750{,}000 \longrightarrow \$350{,}000.
\end{equation}
The mechanism behind the collapse is specific enough to keep: stock options during the dot-com IPO era, insider lockup after the offering, the crash arriving during the lockup, and final liquidation only after severe deterioration.

\subsection{Question \& Answer}

When is wealth only paper wealth?

When it is visible in a quotation but unavailable as cash on the owner's own timetable. In this segment, the peak number is not fictitious, but it is not operationally available either. The lockup transforms price into a constrained observation.

The transcript values let us calculate the scale of the collapse. Using the reported amounts,
\begin{align}
\frac{V_{\text{post-crash}}}{V_{\text{peak}}}
&\approx \frac{0.75}{17.5}
\approx 0.043,\\
\frac{V_{\text{sale}}}{V_{\text{peak}}}
&\approx \frac{0.35}{17.5}
\approx 0.020.
\end{align}
Thus the eventual sale value is roughly \(2\%\) of the paper peak. Equivalently,
\begin{equation}
1-\frac{V_{\text{sale}}}{V_{\text{peak}}} \approx 0.98.
\end{equation}
So about \(98\%\) of the paper peak disappears before realization. This arithmetic is why the speaker's final maxim, ``sell too soon,'' is so sharp. It is not a universal law of asset management. It is a warning that marked wealth under sale constraints is not yet secured wealth.

Only after this loss arithmetic does the lecture return to politics. The speaker declines to discuss politics publicly but says he does not want Kamala Harris to win. The form is consistent: first business mechanism, then political preference.

\section{Election Markets as a Probability Interruption}

The lawyer segment leads directly into the lecture's sponsor-style interruption. The host first repeats the domestic lesson in the broadest possible form---the person one marries may be the most important decision one ever makes---and then abruptly pivots: for the first time, he says, one can legally bet on the election. This interruption should not be mistaken for the chapter's main doctrine, but it matters because it translates political uncertainty into explicit price language.

The stable numerical claims of the segment are these:
\begin{equation}
p_{\text{general}} \approx 50\text{--}50,\qquad
\delta_{\text{race}} \approx 1\%\text{--}2\%,\qquad
V_{\text{platform}} \approx \$1\,\text{billion}.
\end{equation}
The host also gives the promotional deposit-credit arithmetic
\begin{equation}
D=\$100 \Rightarrow C=\$20.
\end{equation}
These are host-side platform claims rather than guest doctrine, but they still mark a genuine change in the lecture's language. We are no longer hearing only about personal experience. We are briefly hearing about public uncertainty as something that can be priced and traded.

The visual evidence for this shift is the dashboard of market cards.

\begin{figure}[t]
\centering
\includegraphics[width=0.95\textwidth]{lecture_61_figure_02.png}
\caption{Prediction-market dashboard with visible state-level presidential odds. The transcript mentions swing states such as Georgia and Arizona, but the readable cards shown here are California, Iowa, New Mexico, Washington, D.C., and Indiana.}
\end{figure}

The screenshot is important not because it proves a theorem, but because it shows the lecture's probabilistic interruption in its native form: not as a map, but as a grid of parallel market cards. The natural notation is
\begin{equation}
p_{s,c}=\Pr(\text{candidate }c\text{ wins state }s).
\end{equation}
Reading only what is clearly visible, we obtain
\begin{align}
p_{\text{CA,Harris}} &\approx 0.98, & p_{\text{CA,Trump}} &\approx 0.02,\\
p_{\text{IA,Trump}} &\approx 0.95, & p_{\text{IA,Harris}} &\approx 0.05,\\
p_{\text{NM,Harris}} &\approx 0.89, & p_{\text{NM,Trump}} &\approx 0.11,\\
p_{\text{DC,Harris}} &\approx 0.98, & p_{\text{DC,Trump}} &\approx 0.02,\\
p_{\text{IN,Trump}} &\approx 0.97, & p_{\text{IN,Harris}} &\approx 0.03.
\end{align}
A side market visible on the same dashboard shows that the platform is not exclusively political:
\begin{equation}
\Pr\!\bigl(\text{RT score of ``Venom: The Last Dance''}>30\bigr)\approx 0.91,
\qquad
\Pr\!\bigl(\text{RT score}>45\bigr)\approx 0.14.
\end{equation}
The small dollar-denominated figures beneath the cards likely refer to some form of market size, liquidity, or traded amount, but the screenshot alone does not settle the exact semantics. What it does settle is the structure: election talk has briefly become a field of explicit probabilities.

This is why the interruption should stay in the chapter. It is the lecture's one clean moment in which belief, uncertainty, and price are put on the page together.

\section{Bankruptcy, Sweat Equity, and the Founder Who Must Be Everyone}

The lecture then returns from public prices to personal rebuilding. The next speaker says he became a millionaire at \(27\), filed bankruptcy at \(31\), and took about ten years to build back up. This gives us the chapter's clearest entrepreneurial timeline:
\begin{align}
a_{\text{startup}} &= 24,\qquad
a_{\text{millionaire}} = 27,\qquad
a_{\text{bankruptcy}} = 31,\\
R_{\text{excavation}} &\approx \$8\,\text{million/year at } a\approx 28\text{--}29,\qquad
T_{\text{rebuild}} \approx 10\ \text{years}.
\end{align}

\paragraph{Anecdote.}
The business was excavation. The speaker says he liked being outside, liked building something, and had tried other trades before discovering that moving earth was the kind of construction he actually wanted to do. The host's line about riches being in the niches lands here because the interview has earned it.

\paragraph{Claim.}
The explicit lesson is that entrepreneurship demands more effort than outsiders think, because there is nobody else above the founder to absorb the work. ``Nobody's coming to save you'' is the host's paraphrase, but it is faithful to the speaker's account of what the role requires.

\paragraph{Mechanism.}
The startup mechanism is sweat equity. The speaker says he worked the equipment by day and went to school at night to learn project management, estimating, and contract negotiation. A cautious compression is
\begin{equation}
K_0^{\text{effective}} \approx K_0^{\text{labor}} + K_0^{\text{skills}},
\end{equation}
with the understanding that this is a note-writer's rendering of the story rather than a formula stated by the speaker himself.

The same idea appears organizationally:
\begin{equation}
\text{founder}
=
\text{owner}
+
\text{VP}
+
\text{operator}
+
\text{estimator}
+
\text{negotiator}.
\end{equation}
This is rhetorical rather than formal, but the rhetoric is exact. Early entrepreneurial competence is multiplex. One is not waiting for the proper department to appear.

\subsection{Question \& Answer}

What survives bankruptcy when the money does not?

The lecture's answer is not optimism. It is specificity. What may survive is operating capacity: habits, learned skills, the ability to estimate, negotiate, manage, and build. The host's recap immediately after the interview says this in nearly theorem-like form. One can lose the business and the money, but not necessarily the skill set and mindset that built them.

We can therefore state the point without overformalizing it:
\begin{equation}
\text{loss of capital does not by itself erase productive capability.}
\end{equation}
The proof supplied by the lecture is biographical. The speaker reaches wealth, loses everything in the crisis, pauses, rebuilds, and reappears not as a blank slate but as someone who still knows how to work. The lecture does not promise recovery for everyone. It distinguishes, however, between the destruction of a balance sheet and the destruction of a builder.

The political answer again follows the business doctrine. The speaker says he is voting for Trump and predicts a landslide. The host then recaps the Austin tally and underscores his surprise that the score so far is uniformly pro-Trump in a city he treats as politically bluer than the surrounding state.

\section{Restaurants, Trading, and the Save-versus-Invest Split}

The last major interview lowers the volume but widens the scope. The speaker is retired from finance, says he was a bond trader, and also says he once owned a small restaurant group.

\paragraph{Anecdote.}
The host asks what kind of business he owned and gets a concrete answer: a couple of restaurants, three locations in all. The transcript then turns quickly toward scale and risk:
\begin{equation}
N_{\text{locations}} = 3,\qquad
R_{\max} \approx \$50\,\text{million}.
\end{equation}
Asked about restaurant margins, the speaker gives a stark two-state description:
\begin{equation}
\text{restaurant outcome} \in \{\text{hugely successful},\ \text{out of business}\}.
\end{equation}
He adds, in that same spirit, that banks do not like lending into restaurants and that the business is privately financed.

\paragraph{Claim.}
Before giving the trading lesson, the speaker makes a useful educational distinction. What one learns in school is important, but what one learns in business is learned hands-on. This is an important transition. The chapter has already passed through real estate, divorce, IPO lockups, and bankruptcy. The final speaker now names the epistemic pattern beneath them: business knowledge is often acquired by doing and losing.

\paragraph{Mechanism.}
That is why the trading maxim lands so cleanly. If a great trader says he never lost money, he is lying. Loss is not an embarrassing side note to learning; it is part of the process by which judgment becomes real. Restaurant risk and trading risk differ in form, but they share a feature: they make downside visible.

\subsection{Question \& Answer}

Why do restaurants and traders both force respect for loss?

Because both place the decision-maker inside a short feedback loop. In restaurants, margins, financing constraints, and operating mistakes expose bad assumptions quickly. In trading, price movement does the same thing. The environments differ, but in each one the world returns an answer before the operator has time to maintain illusions for very long.

That local doctrine prepares the final aphorism. Asked what separates the middle class from wealthy people, the speaker answers:
\begin{equation}
\text{middle class} \to \text{saves},
\qquad
\text{upper class} \to \text{invests}.
\end{equation}
A minimal arithmetic rendering, useful as a reconstruction but not spoken verbatim, is
\begin{align}
W_{t+1}^{\text{save}} &= W_t + s_t,\\
W_{t+1}^{\text{invest}} &= W_t + s_t + r_t W_t.
\end{align}
The difference between the two one-period rules is therefore
\begin{equation}
W_{t+1}^{\text{invest}} - W_{t+1}^{\text{save}} = r_t W_t.
\end{equation}
This is the simplest mathematical statement of the speaker's class distinction. Saving adds the new contribution \(s_t\). Investing adds the same contribution and also exposes existing wealth to return.

The speaker then gives a deliberately unglamorous investment heuristic. Asked what industry people should look at, he answers in effect: whatever you can see people using. He calls it a Peter Lynch model. He lives in a house with five women, everyone has Apple products, everyone wears Lululemon, and so the mystery of investment is reduced to observation. The point is not that these two names are a universal portfolio. The point is that the world of ordinary use can be read as commercial evidence.

The section ends with a politically softer answer than the previous interviews. The speaker describes a split household, says he votes for the party rather than the individual, identifies as Republican, and still refuses certainty about the final result. This softened close is important. The lecture is ending not with a sharpened political claim, but with a broad money distinction.

\section{Summary}

The chapter begins by spending consequences before causes, and it should end by seeing that the lecture has been building one argument across several different domains. The real-estate developer gives us leverage: failure creates hunger, older structures once offered easier lift, and desperation destroys bargaining power. The lawyer gives us realization: quoted wealth under lockup is not yet captured wealth. The sponsor interruption gives us market language: political uncertainty can be expressed as explicit probabilities distributed across many state cards. The excavation entrepreneur gives us recovery: capital can disappear while learned capability remains. The final finance-and-restaurant speaker gives us the quietest but widest compression of all: the middle class saves, the upper class invests.

The official close of the lecture returns to channel promotion and the host's broader program of interviewing wealthy Americans. Conceptually, however, the work is already done. Wealth here is not a static stock. It is a moving sequence of negotiations, constraints, update rules, losses, recoveries, and judgments about what is real and what is only marked on paper. The votes remain part of the format. The deeper subject is still commercial method under pressure.